% 第11章 デリバティブの価格付け(3)(資料11)
\chapter{デリバティブの価格付け (3):多期間二項モデルとプット・コール・パリティ}

本章では,前章の一期間二項モデルを多期間に拡張し,ヨーロピアン・オプションの
価格式である \textbf{CRRの公式}(Cox-Ross-Rubinsteinの公式)を導出する.
さらに,モデルに依存しない重要な関係式である\textbf{プット・コール・パリティ}を
学ぶ.付録として,二項モデルによる信用リスクの評価を紹介する.

\section{多期間二項モデル}

\subsection{モデルの設定}

原資産価格の変動を,二項モデルを多期間つなげた\textbf{二項ツリー}(格子)で
表現する.
\begin{itemize}
\item $p$:上昇確率,$1-p$:下落確率.
\item $u$:上昇ファクター,$d$:下落ファクター.
\item 時点 $t=0$ の価格 $S$ から出発し,各期間で価格は $u$ 倍または $d$ 倍に
なる.例えば $t=3$ では $u^3S,\ u^2dS,\ ud^2S,\ d^3S$ の4状態をとる.
\end{itemize}
$\bigcirc$印の格子(ノード)で状態を表現し,ノードを接続して,時間の進行に伴う
状態変化を示す.

\begin{figure}[htbp]
\centering
\begin{tikzpicture}[scale=0.9,
  nd/.style={draw, circle, minimum size=7mm, inner sep=0.5pt, font=\small}]
\node[nd] (S) at (0,0) {$S$};
\node[nd] (u1) at (2.4,1.2) {$uS$};
\node[nd] (d1) at (2.4,-1.2) {$dS$};
\node[nd] (u2) at (4.8,2.4) {$u^2S$};
\node[nd] (m2) at (4.8,0) {$udS$};
\node[nd] (d2) at (4.8,-2.4) {$d^2S$};
\node[nd] (u3) at (7.2,3.6) {$u^3S$};
\node[nd] (a3) at (7.2,1.2) {$u^2dS$};
\node[nd] (b3) at (7.2,-1.2) {$ud^2S$};
\node[nd] (d3) at (7.2,-3.6) {$d^3S$};
\foreach \a/\b in {S/u1, S/d1, u1/u2, u1/m2, d1/m2, d1/d2,
  u2/u3, u2/a3, m2/a3, m2/b3, d2/b3, d2/d3}{
  \draw[-{Stealth}] (\a) -- (\b);}
\node at (0,-4.4) {\small $t=0$};
\node at (2.4,-4.4) {\small $t=1$};
\node at (4.8,-4.4) {\small $t=2$};
\node at (7.2,-4.4) {\small $t=3$};
\end{tikzpicture}
\caption{多期間二項モデル(3期間の例).各枝の上向きの遷移確率が $p$,
下向きが $1-p$ である.}
\label{fig:binomial-tree}
\end{figure}

\subsection{二期間二項モデル}

まず二期間の場合を考える($t=2$ が満期).原資産のツリーと同じノード上に,
対応する時点・状態における金融商品の価格・キャッシュフローを表すツリーを考える.
満期では,各ノードごとにペイオフが確定している.コールオプションでは,満期の
ペイオフは上から順に $\max(u^2S-K,0)$, $\max(udS-K,0)$, $\max(d^2S-K,0)$ で
ある.

\paragraph{$t=1$ の上昇ノードにおける評価.}
$t=1$ で原資産価格 $uS$ のノード上の複製ポートフォリオを求める.無リスク資産の
利回りを $r$,現在価格は1とする.コールオプションを原資産 $\Delta$ 単位,
無リスク資産 $B$ 単位で複製できたとすると,無裁定条件より次の式が成り立つ:
\begin{align}
\Delta\,u^2S + B(1+r) &= \max(u^2S-K,\,0) \equiv P(uu),\notag\\
\Delta\,udS + B(1+r) &= \max(udS-K,\,0) \equiv P(ud).
\end{align}
この連立方程式を解くと,
\begin{equation}
\Delta = \frac{1}{uS}\,\frac{P(uu)-P(ud)}{u-d},\qquad
B = \frac{1}{1+r}\,\frac{u\,P(ud) - d\,P(uu)}{u-d}
\end{equation}
となり,複製ポートフォリオが得られた.一期間モデルと比較すると,$\Delta$ の
分母の「原資産価格の違い」と分子の「CFの違い」が,このノードのものに変わって
いるだけである.

コールオプションの $t=1$ の上昇ノードにおける価値は,
\begin{align}
C_u = uS\times\Delta + 1\times B
&= \frac{P(uu)-P(ud)}{u-d} + \frac{1}{1+r}\,\frac{uP(ud)-dP(uu)}{u-d}\notag\\
&= \frac{(1+r-d)\,P(uu) - (1+r-u)\,P(ud)}{(u-d)(1+r)}
\end{align}
であり,この価格は次のように書ける:
\begin{equation}
C_u = \frac{\dfrac{1+r-d}{u-d}P(uu)
+ \Bigl(1 - \dfrac{1+r-d}{u-d}\Bigr)P(ud)}{1+r}.
\end{equation}
分子は,$q = (1+r-d)/(u-d)$ を上昇確率としたときの,コールオプションの
(このノードから見た)次期のペイオフの期待値に相当し,分母は無リスク資産の
利回りによる割引に相当する.

\paragraph{$t=1$ の下落ノードと $t=0$ の評価.}
$t=1$ の下落ノードにおける価値 $C_d$ も同様に得られる.さらに $t=0$ における
現在価値を考える.$t=1$ の値 $C_u$, $C_d$ を「満期1期のペイオフ」とみなして,
同じ複製の議論を行う.コールオプションを原資産 $\Delta$ 単位,無リスク資産
$B$ 単位で複製できたとすると,
\begin{align}
\Delta\,uS + B(1+r) &= C_u,\notag\\
\Delta\,dS + B(1+r) &= C_d
\end{align}
より
\begin{equation}
\Delta = \frac{1}{S}\,\frac{C_u - C_d}{u-d},\qquad
B = \frac{1}{1+r}\,\frac{uC_d - dC_u}{u-d}
\end{equation}
となる(複製ポートフォリオが得られた).

\paragraph{$t=0$ の価値.}
コールオプションの $t=0$ のノードにおける価値は,$q = \dfrac{1+r-d}{u-d}$ と
おくと,
\begin{align}
C_0 = \Delta S + B
&= \frac{(1+r-d)C_u - (1+r-u)C_d}{(u-d)(1+r)}
= \frac{q\,C_u + (1-q)\,C_d}{1+r}\notag\\
&= \frac{1}{1+r}\Biggl[
q\,\frac{qP(uu) + (1-q)P(ud)}{1+r}
+ (1-q)\,\frac{qP(ud) + (1-q)P(dd)}{1+r}\Biggr]\notag\\
&= \frac{q^2P(uu) + 2q(1-q)P(ud) + (1-q)^2P(dd)}{(1+r)^2}.
\label{eq:two-period}
\end{align}
分子は確率 $q$ による満期ペイオフの期待値,分母は現在時点までの2期間分の割引で
ある.すなわち,
\begin{equation}
C_0 = \frac{\E^Q[C_2]}{(1+r)^2}
\end{equation}
が成り立つ.ただし $C_2$ はコールオプションの $t=2$ におけるペイオフ,
$\E^Q[\cdot]$ はリスク中立確率 $q$ による期待値である.

\subsection{一般の $n$ 期間への拡張}

前項の結論は次のように一般化できる:
\begin{equation}
C_0 = \frac{\E^Q[C_n]}{(1+r)^n}.
\label{eq:crr-general-call}
\end{equation}
ただし $C_n$ はコールオプションの $t=n$ におけるペイオフである.この式は,
コールに限らず,\textbf{一般の証券に関しても成り立つ}.$t=n$ のペイオフ関数を
原資産価格の関数 $f(S_n)$ とすると,$t=0$ における現在価値 $V_0$ は
\begin{equation}
V_0 = \frac{\E^Q[f(S_n)]}{(1+r)^n}
= \frac{1}{(1+r)^n}\sum_{k=0}^{n}
{}_n\mathrm{C}_k\,q^k(1-q)^{n-k}\,f(u^k d^{n-k}S),
\qquad q = \frac{1+r-d}{u-d}
\label{eq:crr-general}
\end{equation}
で与えられる.ここで ${}_n\mathrm{C}_k\,q^k(1-q)^{n-k}$ は「上昇 $k$ 回,下落
$n-k$ 回のノードに達する確率」である.

\subsection{CRRの公式(コール・オプション)}

再度,コールオプションを考える.ペイオフは
$f(u^kd^{n-k}S) = \max(u^kd^{n-k}S - K,\,0)$ である.ここで,
\begin{quote}
$a$:満期でペイオフ関数が正の値をとるノードのうち,原資産価格が最低のノードを
示す $k$ の値(このノードで,ペイオフがはじめて$+$になる)
\end{quote}
を定義する.$k\geq a$ のノードでは $\max(u^kd^{n-k}S-K,0) = u^kd^{n-k}S-K > 0$,
$k<a$ のノードではペイオフは0である.

このとき,
\begin{align}
V_0 &= \frac{\E^Q[\max(S_n - K,\,0)]}{(1+r)^n}
= \frac{1}{(1+r)^n}\sum_{k=0}^{n}{}_n\mathrm{C}_k\,q^k(1-q)^{n-k}
\max(u^kd^{n-k}S - K,\,0)\notag\\
&= \frac{1}{(1+r)^n}\sum_{k=a}^{n}{}_n\mathrm{C}_k\,q^k(1-q)^{n-k}
(u^kd^{n-k}S - K)\notag\\
&= S\sum_{k=a}^{n}{}_n\mathrm{C}_k\,
\Bigl(\frac{qu}{1+r}\Bigr)^k\Bigl(\frac{(1-q)d}{1+r}\Bigr)^{n-k}
- \frac{K}{(1+r)^n}\sum_{k=a}^{n}{}_n\mathrm{C}_k\,q^k(1-q)^{n-k}.
\label{eq:crr-derivation}
\end{align}

ここで
\begin{equation}
q^* \equiv \frac{qu}{1+r}
\end{equation}
とおくと,
\begin{equation}
1 - q^* = \frac{1+r - qu}{1+r}
= \frac{(1+r)(u-d) - (1+r-d)u}{(u-d)(1+r)}
= \frac{d\,(u - 1 - r)}{(u-d)(1+r)}
= \frac{(1-q)\,d}{1+r}
\end{equation}
が成り立つ(すなわち $q^*$ も確率とみなせる).なお,
\[
q = \frac{1+r-d}{u-d}
\;\Longleftrightarrow\;
qu + (1-q)d = 1+r
\;\Longleftrightarrow\;
\frac{qu}{1+r} + \frac{(1-q)d}{1+r} = 1
\]
であり,真ん中の式は「(確率 $q$ のもとで)原資産の期待収益率 $= r$」,つまり
\textbf{リスク中立確率を使っている}という意味である.

さらに,二項モデルの\textbf{補分布関数}(離散分布における生存関数)
\begin{equation}
\phi(a; n, p) \equiv \sum_{k=a}^{n}{}_n\mathrm{C}_k\,p^k(1-p)^{n-k}
\end{equation}
を定義する.すると,\eqref{eq:crr-derivation}から
\begin{equation}
\boxed{\;
V_0 = S\,\phi(a; n, q^*) - \frac{K}{(1+r)^n}\,\phi(a; n, q)
\;}
\label{eq:crr-call}
\end{equation}
が得られる.ただし,$a$ は $u^ad^{n-a}S > K$ を満たす最小の整数であり,
\begin{equation}
\Bigl(\frac{u}{d}\Bigr)^a > \frac{K}{d^nS}
\;\Longleftrightarrow\;
a\log\frac{u}{d} > \log\frac{K}{d^nS}
\qquad\therefore\quad
a > \frac{\log(K/S) - n\log d}{\log(u/d)}
\label{eq:crr-a}
\end{equation}
を満たす最小の整数が $a$ である.

\subsection{CRRの公式(プット・オプション)}

同じ多期間二項モデルで,ヨーロピアン・プット・オプションの価格式(CRRの公式の
プットオプション版)を求める.以下,コールオプションの価格式導出と対応させて
書く(講義では自習箇所とされた).

一般化された結果
\begin{equation}
P_0 = \frac{\E^Q[P_n]}{(1+r)^n},\qquad
V_0 = \frac{\E^Q[f(S_n)]}{(1+r)^n}
= \frac{1}{(1+r)^n}\sum_{k=0}^{n}{}_n\mathrm{C}_k q^k(1-q)^{n-k}f(u^kd^{n-k}S)
\end{equation}
を使う($P_n$ はプットオプションの $t=n$ におけるペイオフ).ペイオフ関数は
$f(u^kd^{n-k}S) = \max(K - u^kd^{n-k}S,\,0)$ であり,プットのペイオフが正になる
のは $k \leq a-1$ のノード($a$ はコールと同じ,$u^ad^{n-a}S>K$ を満たす最小の
整数)である.このとき,
\begin{align}
V_0 &= \frac{\E^Q[\max(K-S_n,\,0)]}{(1+r)^n}
= \frac{1}{(1+r)^n}\sum_{k=0}^{a-1}{}_n\mathrm{C}_k\,q^k(1-q)^{n-k}
(K - u^kd^{n-k}S)\notag\\
&= \frac{K}{(1+r)^n}\sum_{k=0}^{a-1}{}_n\mathrm{C}_k\,q^k(1-q)^{n-k}
- S\sum_{k=0}^{a-1}{}_n\mathrm{C}_k\,
\Bigl(\frac{qu}{1+r}\Bigr)^k\Bigl(\frac{(1-q)d}{1+r}\Bigr)^{n-k}.
\end{align}
ここで,離散分布における分布関数
\begin{equation}
F(a; n, p) \equiv \sum_{k=0}^{a-1}{}_n\mathrm{C}_k\,p^k(1-p)^{n-k}
= 1 - \phi(a; n, p)
\end{equation}
を定義すると,
\begin{equation}
\boxed{\;
P_0 = \frac{K}{(1+r)^n}\,F(a; n, q) - S\,F(a; n, q^*)
\;}
\label{eq:crr-put}
\end{equation}
が得られる.

\begin{definition}[CRRモデルとCRRの公式]
ここまでで示してきた,多期間二項モデルによるオプションの価格付けモデルを
\textbf{CRRモデル}(Cox-Ross-Rubinsteinモデル)という.特に,彼らが導出した
ヨーロピアン・コール(プット)・オプション価格式\eqref{eq:crr-call},
\eqref{eq:crr-put}を \textbf{CRRの公式}(Cox-Ross-Rubinsteinの公式)という.
オプション価格の一般式\eqref{eq:crr-general}もCRRの公式と呼ぶことがある
(どこまでを「CRRの公式」と呼ぶかは明確でない).
\end{definition}

\section{補足:CRR公式の正確な表現(重要!)}

金利による割引の表現に注意が必要である.これまでは一期間の利子率を $r$ として
表示してきたが,本来,一期間を $\Delta t$ 年,満期を $n\Delta t$ 年,金利を
\emph{年率} $r$ として,
\begin{equation}
C_0 = \frac{\E^Q[C_n]}{(1+r\Delta t)^n}
\end{equation}
が正確な表現である(これまでの $r$ を $r\Delta t$ に置換する).同様にして,
一般化されたCRRの価格式は
\begin{equation}
V_0 = \frac{\E^Q[f(S_n)]}{(1+r\Delta t)^n}
= \frac{1}{(1+r\Delta t)^n}\sum_{k=0}^{n}
{}_n\mathrm{C}_k\,q^k(1-q)^{n-k}f(u^kd^{n-k}S),
\qquad
q = \frac{1 + r\Delta t - d}{u - d}
\end{equation}
が正確な表現である.

\paragraph{割引金利が一定でない場合.}
さらに,第 $i$ 期目($t_{i-1}$ から $t_i$ まで)の無リスク金利が確定的で,
年率 $r_i$ で与えられる(時点によって異なる)場合は,
\begin{equation}
C_0 = \frac{\E^Q[C_n]}{\prod_{i=1}^{n}(1 + r_i\Delta t)}
\end{equation}
になる.\emph{しかし},CRRの価格式は
\[
V_0 = \frac{\E^Q[f(S_n)]}{\prod_{i=1}^n(1+r_i\Delta t)}
= \frac{1}{\prod_{i=1}^n(1+r_i\Delta t)}
\sum_{k=0}^{n}{}_n\mathrm{C}_k\,q^k(1-q)^{n-k}f(u^kd^{n-k}S)
\]
とは\textbf{ならない}.一般に,確率 $q$ が無リスク金利 $r_i$ に依存して変わる
ため,確率分布を二項分布では表現できなくなるからである.

\section{プット・コール・パリティ}

\subsection{プットとコールのペイオフの関係}

コール・ロングとプット・ロングのペイオフを比べると,原資産,満期,行使価格が
等しいならば,
\[
\text{コール・ロングのペイオフ} - \text{プット・ロングのペイオフ}
= \text{原資産価格} - K,
\]
すなわち
\begin{equation}
\max(S(T)-K,\,0) - \max(K-S(T),\,0) = S(T) - K
\label{eq:payoff-identity}
\end{equation}
が成り立つ(もともとこれは恒等式である.$S(T)\geq K$ の場合と $S(T)<K$ の場合に
分けて確かめればよい).

\subsection{プットオプションの複製}

式\eqref{eq:payoff-identity}を変形した
\begin{equation}
\max(K - S(T),\,0) = \max(S(T)-K,\,0) - S(T) + K
\end{equation}
の左辺をプットのペイオフ,右辺を複製ポートフォリオのペイオフとみると,
「プットを1単位ロング」の複製ポートフォリオが
\begin{itemize}
\item[A.] コールを1単位ロング,
\item[B.] 原資産を1単位ショート,
\item[C.] 満期 $T$ の割引債を $K$ 単位ロング,
\end{itemize}
であることを示している.Cで,将来のキャッシュを「割引債」のロング(ケースに
よってはショート)で表現している点に注意する.

\subsection{プット・コール・パリティ}

無裁定の議論より,「オプションとその複製ポートフォリオは,現時点でも価格は
等しい」ので,次式が成り立つ.

\begin{theorem}[プット・コール・パリティ]
\begin{equation}
p(t) = c(t) - S(t) + K\,v(t,T),
\label{eq:put-call-parity}
\end{equation}
あるいは書き直して
\begin{equation}
p(t) + S(t) = c(t) + K\,v(t,T).
\end{equation}
ここで,$p(t)$:時刻 $t$ におけるヨーロピアンプットオプションの価格,
$c(t)$:時刻 $t$ におけるヨーロピアンコールオプションの価格,
$S(t)$:時刻 $t$ における原資産価格,
$v(t,T)$:時刻 $t$ における満期 $T$ の(無リスクな)割引債価格である.
\end{theorem}

この関係式を\textbf{プット・コール・パリティ}という.これは価格付けモデルに
依存しない,一般に成り立つ関係式である.コール価格が得られれば,プット価格は
この式から得られる.逆も同様である.

\section*{付録:信用リスクの評価(参考)}

二項モデルによる信用リスクの評価を紹介する(講義では扱わない参考資料である).

\paragraph{信用リスクとは.}
\begin{itemize}
\item \textbf{デフォルト}:発行体や取引先が債務を履行できなくなる状態.
\item \textbf{信用リスク} (credit risk):デフォルトにより被る直接的・間接的な
損失.
  \begin{itemize}
  \item 直接的な損失:将来の契約が履行されなくなることによる損失.
  \item 間接的な損失:デフォルトの可能性の高まりにより,その証券の取引価格が
  下落すること.
  \end{itemize}
\end{itemize}

デフォルトによる直接的な損失の例として,クーポン3円・満期6年・額面100円の債券を
考えると,予定キャッシュフローは毎年3円+満期に103円だが,例えば4年目に
デフォルトが発生すると,極端な場合,デフォルト発生以降のキャッシュフローはゼロに
なる.

\paragraph{例:割引社債の二項モデル評価.}
一般事業会社Aが満期1年の割引社債を発行したとする.投資家は,デフォルトしなければ
1年後に1円受け取り,満期前にデフォルトしたら1年後に $\delta$ 円受け取ると仮定
する($0\leq\delta\leq 1$.$\delta$ は\textbf{回収率}と呼ばれる).

この割引社債の評価を二項モデルで考える.二項モデルの将来ノードは,ペイオフに
関連する状態を表していればよい.複数のノードでペイオフが等しくなってもよいし,
必ずしも原資産価格の水準の変化だけで考えなくてよい.柔軟に考えるほど,応用先は
広がる.

(リスク中立確率下の)デフォルト確率を $q$,デフォルト時の回収率を $\delta$ と
する.$t=1$ で,(1) 常に1円受け取る証券と,(2) デフォルト時のみ1円受け取る
証券を考える.(1) は無リスクな債券(割引債,現在価格 $1/(1+r)$),(2) は
CDS (Credit Default Swap) である.「デフォルト時にのみ1円受け取れる証券
(CDS)」が存在すれば,そのCDSと無リスク証券で社債を複製可能である.

無リスク金利 $r$ を一定とすると,リスク中立化法より,社債の現在価格は
\begin{equation}
\pi = \frac{\E^Q[C]}{1+r} = \frac{(1-q)\times 1 + q\times\delta}{1+r}.
\end{equation}
この式は
\begin{equation}
\pi = \frac{1}{1+r} - (1-\delta)\,\frac{q}{1+r}
\end{equation}
と書ける.第1式は「キャッシュフローの期待値の割引価値が現在価格になる」ことを
示し,第2式は複製ポートフォリオを示している:第1項は無リスクの割引債(これも
無リスク資産の一表現)を1単位ロング,第2項はCDSを $(1-\delta)$ 単位ショート,
である.

% ---------------- 演習問題 ----------------
\section{演習問題}

\begin{exercise}[問題11-1]
多期間二項モデルのCRRの公式で,プット・コール・パリティを確認しなさい.
\end{exercise}

\begin{answerbox}
\textbf{考え方.}CRRのプット価格式\eqref{eq:crr-put}に
$F(a;n,p) = 1-\phi(a;n,p)$ を代入し,プット・コール・パリティの右辺
$c - S + K(1+r\Delta t)^{-n}$ の形になることを計算で確かめる.

\medskip
\textbf{解答.}
一期間を $\Delta t$ 年とする正確な表現で書く.CRRの公式より
\begin{align*}
\mathrm{Put} + S - \frac{K}{(1+r\Delta t)^n}
&= \frac{K}{(1+r\Delta t)^n}F(a;n,q) - S\,F(a;n,q^*)
+ S - \frac{K}{(1+r\Delta t)^n}\\
&= \frac{K}{(1+r\Delta t)^n}\bigl(F(a;n,q) - 1\bigr)
+ S\bigl(1 - F(a;n,q^*)\bigr)\\
&= -\frac{K}{(1+r\Delta t)^n}\,\phi(a;n,q) + S\,\phi(a;n,q^*)\\
&= \mathrm{Call}.
\end{align*}
すなわち
\[
\mathrm{Put} = \mathrm{Call} - S + \frac{K}{(1+r\Delta t)^n}
\]
が成り立つ.右辺の第3項 $K(1+r\Delta t)^{-n}$ は満期 $n\Delta t$ の割引債
$K$ 単位の現在価値 $K\,v(0,T)$ にほかならないから,これはプット・コール・
パリティ\eqref{eq:put-call-parity}そのものである.$\blacksquare$

\medskip
\textbf{ポイント.}CRRの公式は特定のモデル(二項モデル)による価格式だが,
そこから得られるコール価格とプット価格は,モデルフリーな関係式である
プット・コール・パリティをきちんと満たしている.価格モデルの整合性チェックの
好例である.
\end{answerbox}

\begin{exercise}[問題11-2]
ここで説明した多期間二項モデルを使い,満期 $T$,バリアー価格 $B$ で $A$ 円を
受け取るデジタル・オプションの価格式を求めなさい.

デジタル・オプションとは,満期 $T$ における原資産価格 $S(T)$ がバリアー価格
$B$ を上回るとき,満期 $T$ に一定額 $A$ 円を受け取る証券である.$S(T)<B$ の
ときは何も受け取れない.
\end{exercise}

\begin{answerbox}
\textbf{考え方.}CRRの一般式\eqref{eq:crr-general}に,デジタル・オプションの
ペイオフ関数を代入すればよい.ペイオフはノードの原資産価格が $B$ を上回るか
どうかだけで決まるので,和はペイオフが正になるノードだけに絞られる.

\medskip
\textbf{解答.}
満期 $T$ におけるキャッシュフローは
$f(S(T)) = A\,\bm{1}_{\{S(T)>B\}}$ なので,CRRの公式より
\begin{align*}
V_0 &= \frac{\E^Q[f(S_n)]}{(1+r\Delta t)^n}
= \frac{1}{(1+r\Delta t)^n}\sum_{k=0}^{n}
{}_n\mathrm{C}_k\,q^k(1-q)^{n-k}\,A\times\bm{1}_{\{u^kd^{n-k}S > B\}}\\
&= \frac{A}{(1+r\Delta t)^n}\sum_{k=b}^{n}{}_n\mathrm{C}_k\,q^k(1-q)^{n-k}
= \frac{A}{(1+r\Delta t)^n}\,\phi(b;n,q).
\end{align*}
(注:$\bm{1}_A$ は定義関数で,事象 $A$ が真のとき1,偽のとき0.)
ただし,$b$ は $u^bd^{n-b}S > B$ を満たす最小の整数で,
\[
\Bigl(\frac{u}{d}\Bigr)^b > \frac{B}{d^nS}
\;\Longleftrightarrow\;
b\log\frac{u}{d} > \log\frac{B}{d^nS}
\qquad\therefore\quad
b > \frac{\log(B/S) - n\log d}{\log(u/d)}.
\]
($B\to K$ とみなすと,$b$ はこれまで使ってきた $a$ と同じである.)

\medskip
\textbf{ポイント.}コールオプションの価格式\eqref{eq:crr-call}の第2項
(行使価格の支払い部分)は,実は「$S(T)>K$ のとき $K$ 円支払う」デジタル・
オプションの価値である.デジタル・オプションはCRR公式の構成要素そのものと
言える.
\end{answerbox}

\begin{exercise}[問題11-3]
ここで説明したCRRモデルを用いて,
\begin{itemize}
\item 原資産価格 $S(0) = 100$ 円,
\item 満期1年,
\item 行使価格 $K = 100$ 円,
\item $\Delta t = 1/4$ 年($=$ 3ヶ月),
\item 無リスク金利 4\%(年率),
\item $u = 1.1$, $d = 0.9$
\end{itemize}
のときのヨーロピアン・コール・オプションの価格を求めなさい.
\end{exercise}

\begin{answerbox}
\textbf{考え方.}$n = 1/\Delta t \times 1$年 $= 4$ 期間の二項モデルである.
CRRの公式(正確な表現)を使うため,(1) リスク中立確率 $q$,(2) $q^*$,
(3) ペイオフが正になる最小ノード $a$,を順に計算し,公式\eqref{eq:crr-call}に
代入する.

\medskip
\textbf{解答.}
まず,原資産価格のリスク中立的な上昇確率は
\[
q = \frac{1 + r\Delta t - d}{u - d}
= \frac{1 + 0.04\times 0.25 - 0.9}{1.1 - 0.9} = \frac{0.11}{0.2} = 0.55
\]
より,55\%である.さらに,
\[
q^* = \frac{qu}{1+r\Delta t} = \frac{0.55\times 1.1}{1 + 0.04\times 0.25}
= \frac{0.605}{1.01} = 0.59901\ldots.
\]
次に,ペイオフが正になるノードを示す $k=a$ は,
\[
a > \frac{\log(K/S) - n\log d}{\log(u/d)}
= \frac{\log(100/100) - 4\log 0.9}{\log(1.1/0.9)} = 2.100\ldots
\]
より,$a = 3$ である.

CRRの公式より,コール・オプション価格は
\begin{align*}
V_0 &= S\,\phi(a;n,q^*) - \frac{K}{(1+r\Delta t)^n}\,\phi(a;n,q)\\
&= 100\,\phi(3;4,\ 0.59901) - \frac{100}{(1+0.04\times 0.25)^4}\,
\phi(3;4,\ 0.55)\\
&= 100\times 0.473490 - \frac{100}{1.01^4}\times 0.390981\\
&\approx 9.7765 \text{ 円}.
\end{align*}
ここで
\[
\phi(3;4,q^*) = {}_4\mathrm{C}_3\,(q^*)^3(1-q^*) + (q^*)^4 = 0.473490,
\qquad
\phi(3;4,q) = {}_4\mathrm{C}_3\,q^3(1-q) + q^4 = 0.390981
\]
である.

\medskip
\textbf{検算(4期間二項ツリーによる後ろ向き帰納法).}
満期のペイオフは,上から
$u^4S - K = 46.41$,\ $u^3dS - K = 19.79$,\ 0,\ 0,\ 0 であり,各ノードで
$V = [qV_u + (1-q)V_d]/(1+r\Delta t)$ を繰り返すと,
\[
t=3:\ (34.0901,\ 10.77673,\ 0,\ 0),\quad
t=2:\ (23.36543,\ 5.868518,\ 0),\quad
t=1:\ (15.33844,\ 3.195728)
\]
を経て $t=0$ で $9.776452$ 円となり,CRRの公式による値と一致する.
\end{answerbox}
